\begin{note}
	Отметим, что если элементы ряда будут из $\C$. Тогда теоремы: необходимое условие Коши, критерий Коши, абсолютная сходимость $\ra$ сходимость, признак Вейерштрасса, признаки Дирихле и Абеля справедливы, когда монотонные последовательности состоят из действительных чисел (или функций).
\end{note}
\subsection{Степенные ряды}
\begin{definition}
	\text{Степенным рядом} называется $\sssc$, где коэффиценты $\{c_n\}_{n=0}^\i\cup\{z_0\}\subset\C,\ z\in\C$.
\end{definition}
\begin{definition}
	\textit{Радиусом сходимости} степенного ряда называется:
	\[R=\frac1{\vlim[n\ra\i]\sqrt[n]{|c_n|}}\in[0,+\i]\]
	при соглашении, что $\ds\frac10=+\i,\ \frac1{+\i}=0$.
\end{definition}
\begin{theorem}\nf{(Коши - Адамара)}\label{17_th1}\\
	Степенной ряд $\sssc$:
	\begin{enumerate}
		\item Если $|z-z_0|<R$, где $R$ - радиус сходимости ряда, то ряд сходится абсолютно.
		\item Если $|z-z_0|>R$, то ряд расходится.
	\end{enumerate}
\end{theorem}
\begin{proof}
	\begin{enumerate}
		\item По признаку Коши для ряда $|c_n|\cdot|z-z_0|^n$:
		\[\vlim[n\ra\i]\sqrt[n]{|c_n|\cdot|z-z_0|^n}=|z-z_0|\cdot\vlim[n\ra\i]\sqrt[n]{|c_n|}<1\]
		Значит ряд сходится абсолютно. Также отметим, если $R=+\i$, то $\vlim[n\ra\i]\sqrt[n]{|c_n|}=0$, тогде неравенство тоже выполняется.
		\item Имеем:
		\[\vlim[n\ra\i]\sqrt[n]{|c_n|\cdot|z-z_0|^n}>1\]
		Теперь вспомним доказательство признака Коши. Тогда:
		\[|c_n\cdot(z-z_0)^n|\nra0\]
	\end{enumerate}
	Получили требуемое.
\end{proof}
\begin{corollary}\nf{(Первая теорема Абеля)}\\
	Если степенной ряд $\sssc$ сходится при некотором $z_1\in\C$, то он сходится $\fa z$ такого, что $|z-z_0|<|z_1-z_0|$.
\end{corollary}
\begin{proof}
	Воспользуемся только что доказанной теоремой:
	\[|z-z_0|<|z_1-z_0|\le R\]
	Получили требуемое.
\end{proof}
\begin{theorem}\nf{(Другая формула для радиуса сходимости)}\label{17_th2}\\
	Если $\ds\ex\lim_{n\ra\i}\l|\frac{c_n}{c_{n+1}}\r|=L\in[0,+\i]$, то радиус сходимости ряда равен $L$.
\end{theorem}
\begin{proof}
	Возьмём последовательность $\ds\al_n=\frac{c_{n+1}}{c_n}$.\\
	Тогда по следствию из регулярности суммирования Чезаро знаем:
	\[\lim_{n\ra\i}\l|\al_n\r|=\frac1L=\lim_{n\ra\i}\sqrt[n]{|a_0\dots a_n|}=\lim_{n\ra\i}\sqrt[n]{\frac{c_1}{c_0}\cdot\frac{c_2}{c_1}\cdot\dots\cdot\frac{c_{n+1}}{c_n}}=\lim_{n\ra\i}\sqrt[n]{c_{n+1}}=\lim_{n\ra\i}\sqrt[n]{c_n}\]
	\[R=\frac{1}{1/L}=L\]
	Получили требуемое.
\end{proof}
\begin{theorem}\nf{(Равномерная сходимость степенных рядов)}\label{17_th3}\\
	Если радиус сходимости степенного ряда $\sssc$ $R>0$, то ряд сходится равномерно на круге $|z-z_0|<r$ для любого $0<r<R$.
\end{theorem}
\begin{proof}
	Положим $z_1:=z_0+r$, тогда $|z_1-z_0|=r<R\Ra$ ряд сходится абсолютно в точке $z_1$, то есть:
	\[|c_n(z-z_0)^n|\le|c_n|\cdot r^n\]
	Тогда по признаку Вейерштрасса сходится равномерно на $|z-z_0|\le r$.\\
	Получили требуемое.
\end{proof}
\begin{corollary}
	Сумма ряда непрерывна в каждой точке круга сходимости $|z-z_0|<R$.
\end{corollary}
\begin{proof}
	Все члены ряда непрерывны для $\fa r<R$, тогда непрерывны они и при $|z-z_0|<R$.
\end{proof}
\begin{theorem}\nf{(Вторая теорема Абеля)}\label{17_th4}\\
	Если степенной ряд $\sssc$ сходится при некотором $z_1\in\C$, то его сумма непрерывна на отрезке $[z_0,z_1]$.
\end{theorem}
\begin{proof}
	\[z\in[z_0,z_1]\lra z=z_0+t\cdot(z_1-z_0),\ t\in[0,1]\]
	\[\sssc=\ss[0]c_n(z_1-z_0)^n\cdot t^n\]
	$\ss[0]c_n(z_1-z_0)^n$ равномерно сходится относительно $t,\ \{t^n\}_{n=0}^\i$ убывает, равномерно ограничена.\\
	Тогда, по признаку Дирихле, ряд сходится равномерно, а из этого следует, что сумма непрерывна.\\
	Получили требуемое.
\end{proof}
\begin{corollary}
	Если ряды $\ss[0]a_n,\ \ss[0] b_n$ и их произведение по Коши $\ss[0] c_n$ сходятся, то:
	\[\ss[0] c_n=\l(\ss[0]a_n\r)\cdot\l(\ss[0]b_n\r)\]
\end{corollary}
\begin{proof}
	Рассмотрим степенные ряды:
	\[\ss[0]a_n\cdot z^n,\ \ss[0]b_n\cdot z^n,\ \ss[0]c_n\cdot z^n\]
	Заметим, что произведение по Коши первых 2 совпадёт с третьим:
	\[\ps[0]a_k\cdot z^k\cdot b_{n-k}\cdot z^{n-k}=\l(\ps[0]a_k\cdot b_{n-k}\r)\cdot z^n=c_n\cdot z^n\]
	При $z\in[0,1)$, по (\ref{17_th1}), все эти степенные ряды сходятся абсолютно.\\
	Тогда по теореме о произведении абсолютно сходящихся рядов:
	\[\l(\ss[0]a_n\cdot z^n\r)\cdot\l(\ss[0] b_n\cdot z^n\r)=\ss[0]c_n\cdot z^n\text{ при }z\in[0,1)\]
	Так как функции непрерывны, то выражение выше справедливо и при $z=1$.\\
	Получили требуемое.
\end{proof}
\begin{definition}
	Назовём ядром Дирихле следующее равенство:
	\[\frac12+\ps[1]\cos(k\phi)=
	\frac12+\frac1{2\sin(\frac\phi2)}\cdot
	\ps[1]\l(\sin\l(\frac{2k+1}{2}\phi\r)-\sin\l(\frac{2k-1}{2}\phi\r)\r)
	=\frac{\sin(\frac{2n+1}2\phi)}{2\sin\l(\frac\phi2\r)}\]
	Вспомнить выведение этого равенства можно, если домножить $\cos$ под суммой на $\ds\frac{\sin(\frac\phi2)}{\sin(\frac\phi2)}$ и применить формулу на произведение $\sin$ и $\cos$.
\end{definition}
\begin{examples}
	Покажем, как себя может вести степенной ряд на окружности радиуса сходимости.
	\begin{enumerate}
		\item $\ss[1]\frac{z^n}n,\ R=1$. При $z=1$ ряд расходится.\\
		При $z\neq1,\ |z|=1$ - сходится.\\
		$z=\cos\phi+i\cdot\sin\phi,\ z^n=\cos(n\phi)+i\cdot\sin(n\phi)$
		\[\ss\frac{z^n}n=\ss\frac{\cos(n\phi)}n+i\cdot\ss\frac{\sin(n\phi)}n=a+i\cdot b\]
		Знаем, что $a$ сходится при $\phi\in\R$, а ряд $b$ - при $\phi\in(0,2\pi)$.
		\[(\ex p\in \R)\ \ps[1]\cos(k\phi)=\frac{\sin(\frac{2n+1}{2}\phi)-\sin(\frac\phi2)}{2\sin(\frac\phi2)}\le p\]
		Тогда можем применить признак Дирихле, так как $\frac1n$ убывает к 0.
		Получили требуемое.
		\item $\ds\ss\frac{z^n}{n^2},\ R=1, \l|\frac{z^n}{n^2}\r|\le\frac1{n^2}$.\\
		По признаку Вейерштрасса сходится при $\fa z,\ |z|=1$.
		\item $\ds\ss[1]z^n,\ R=1,\ \fa z,\ |z|=1$. Расходится.
	\end{enumerate}
\end{examples}
\subsection{Действительные степенные ряды}
\begin{note}
	С этого момента рассматриваться будут только степенные ряды с коэффицентами из $\R$:
	\[\sss,\ \{a_n\}_{n=0}^\i\cup\{x_0\}\subset\R,\ x\in\R\]
\end{note}
\begin{lemma}
	Радиусы сходимости степенных рядов $\sss$ и $\ss na_n(x-x_0)^n$ совпадают.
\end{lemma}
\begin{proof}
	\[\vlim[n\ra\i]\sqrt[n]{|na_n|}=\lim_{n\ra\i}\sqrt[n]{n}\cdot\vlim[n\ra\i]\sqrt[n]{|a_n|}\]
	Получили требуемое.
\end{proof}
\begin{theorem}\nf{(Дифференцирование и интегрирование степенных рядов)}\label{17_th5}\\
	Если радиус сходимости степенного ряда $\sss$ равен $R>0$, то справедливы следующие утверждения:
	\begin{enumerate}
		\item $\fa x\in(x_0-R,x_0+R)$ сумма ряда бесконечно дифференцируема, причём можно дифференцировать почленно, то есть, если $f(x)=\sss$, то \[f^{(k)}(x)=\ss[k]a_n\cdot n\cdot (n-1)\cdot\dots\cdot(n-k+1)\cdot(x-x_0)^{n-k}\]
		\item $\fa x\in(x_0-R,x_0+R)$ ряд почленно интегрируем на отрезке с концами $x_0$ и $x$:
		\[\int_{x_0}^xf(t)dt=\ss[0] a_n\cdot(x-x_0)^{n+1}\cdot\frac1{n+1}\]
		\item Все ряды в пунктах 1 и 2 имеют тот же радиус сходимости.
		\item Коэффиценты ряда вычисляются с помощью формулы:
		\[a_n=\frac{f^{(n)}(x_0)}{n!},\ n=0,1,2,\dots\]
	\end{enumerate}
\end{theorem}
\begin{proof}
	Начнём с пункта 3. Почленно продифференцированный ряд имеет вид:
	\[\ss[1]n\cdot a_n(x-x_0)^{n-1}=\frac1{x-x_0}\cdot\ss[1]n\cdot a_n(x-x_0)^n\]
	Радиусы сходимости совпадают. Один раз можно дифференцировать, значит продифференцировав ещё раз получим тот же результат, и так далее.\\
	Для интегрируемости поменяем местами ряды. Получили, что радиусы сходимости совпадают.\\
	Теперь разберёмся со 2 пунктом. По теореме (\ref{16_th7'}) получим требуемое.\\
	1 пункт следует из теоремы (\ref{16_th8'}).\\
	4 пункт:
	\[f^{(k)}(x_0)=a_k\cdot k!\]
	Получили требуемое.
\end{proof}